\section{线程束、处理块与单指令多线程执行}

\subsection{线程块进一步划分为线程束}

线程块驻留到流式多处理器后，硬件按照线程的\textbf{线性化索引}把它拆分成若干\term{线程束}{warp}。线程束大小由设备实现决定；截至 Hopper 的所列 NVIDIA 架构均采用 32 个线程。

若令线程束大小为 \(W=32\)，线程块中线性化索引为 \(t\) 的线程属于
\[
\text{warpId}=\left\lfloor\frac{t}{W}\right\rfloor.
\]
因此，索引 \(0\sim31\) 属于线程束 0，索引 \(32\sim63\) 属于线程束 1，依此类推。一个包含 \(T_b\) 个线程的线程块需要
\[
W_b=\left\lceil\frac{T_b}{W}\right\rceil
\]
个线程束。

\begin{examplebox}[title=例题：三个 256 线程的线程块]
每个线程块包含
\[
256/32=8
\]
个线程束。若三个这样的线程块同时驻留在同一个流式多处理器上，则共有
\[
8\times 3=24
\]
个驻留线程束。这里的 24 表示可供调度的线程束状态数量，不表示 24 个线程束在同一时钟周期同时发射指令。
\end{examplebox}

\subsection{处理块与线程束调度器}

流式多处理器内部的执行核心还会组织成若干\term{处理块}{processing block, PB}。在简化的流式多处理器模型中，每个处理块具有自己的控制与线程束调度资源；线程束被调度到这些处理块上执行。

\begin{figure}[H]
  \centering
  \includegraphics[width=0.92\textwidth]{sm_scheduling_p14.png}
  \caption{线程块先拆分成线程束；流式多处理器内部再由多个处理块及其控制资源调度这些线程束}
\end{figure}

以 Hopper H100 为例，128 个 FP32 核心组织成 4 个处理块，每个处理块包含 32 个 FP32 核心并共享相应的指令分派资源。该图是吞吐结构的简化表示，因为一个流式多处理器内还存在整数单元、加载/存储单元、FP64 单元与张量核心等其他执行资源。

\subsection{从单程序多数据到单指令多线程}

软件层次的\term{单程序多数据}{Single Program, Multiple Data, SPMD}描述的是：大量线程运行同一个核函数程序，每个线程依据自己的索引和数据完成不同实例的工作。

硬件调度层次则使用\term{单指令多线程}{Single Instruction, Multiple Threads, SIMT}：当某个线程束被选中时，调度器为该线程束发射一条指令，各个活动线程在各自的数据和线程状态上执行这条指令。该组织与\term{单指令多数据}{Single Instruction, Multiple Data, SIMD}具有相同的控制复用思想，即让多个执行通道共享指令获取和分派成本。

\begin{keybox}
SPMD 是程序员看到的核函数编程模型；SIMT 是线程束层次的硬件执行模型。二者的连接点是线程束：来自同一核函数的线程被成组调度，同一时刻的指令发射以线程束为单位。
\end{keybox}

这种结构的主要硬件收益是摊薄控制成本。若 32 个执行通道都要执行同一条指令，就不需要为每个通道复制完全独立的指令获取与分派逻辑。代价则是：如果线程束内部的线程需要走不同控制路径，执行通道不能始终保持同时有效，这就产生控制发散。

\subsection{独立线程调度不改变性能分析的基本单位}

从 Volta 开始，NVIDIA 引入\term{独立线程调度}{independent thread scheduling}。线程束中的线程可以维护更独立的执行状态，在发散后沿不同路径前进并在之后重新会合，不再要求所有线程在所有情况下表现为严格锁步。

性能分析仍把线程束视为关键单位，原因有两点：第一，执行资源和指令发射仍围绕线程束组织；第二，线程束内部的控制不一致依然会降低有效通道利用率。独立线程调度还意味着某些过去依赖隐含锁步行为的线程间协作必须使用更明确的同步机制。

\subsection{驻留线程束、就绪线程束与实际发射线程束}

对一个流式多处理器中的线程束，需要区分三个状态概念：

\begin{itemize}
  \item \textbf{驻留线程束}：其线程状态、寄存器等资源已经分配在流式多处理器中；
  \item \textbf{就绪或可选线程束}：下一条指令所需操作数与执行条件已经满足，可以被调度器选择；
  \item \textbf{实际发射线程束}：当前获得发射机会、正在向执行资源送出指令的线程束。
\end{itemize}

占用率统计的是驻留规模，而不是“这一周期有多少线程束正在发射”。理解这一差别，是后续延迟隐藏和占用率分析的前提。
